% Sources: README "Visibility and spectrum"
% Code: ULXspectra, branch fix-results-integrity, commit 5f00569;
%       wind spectrum (eq. Inu_wind): ULXspectra-mc, branch mc-transfer-corrections, commit f546151

\section{Visibility and the observed spectrum}
\label{sec:spectrum}

The emitting surface consists of rings: the wind photosphere and the outer wall, the directly seen disc and, for a NS, the magnetosphere.
For the observer at the inclination $i$, each ring $j$ has a projected area $A_{{\rm proj},j}(i)$.
The surface is axisymmetric, $z(R)$, and we divide each ring into \code{nphi\_area} azimuthal elements.
An element at azimuth $\phi$ contributes its area $R\,{\rm d}R\,{\rm d}\phi$ times the projection of its normal onto the line of sight,
\begin{equation}
\mu_{\rm proj} = \cos i - \frac{{\rm d}z}{{\rm d}R}\,\sin i\,\cos\phi ,
\label{eq:muproj}
\end{equation}
% code: projected_area_weights_surface_raycast() (mu_proj) (csrcs/visibility.c)
if it faces the observer ($\mu_{\rm proj} > 0$) and if the ray from it towards the observer is not blocked by the rest of the surface, including the disc plane beyond $\rout$.
% code: projected_area_weights_surface_raycast(), is_occulted_by_surface() (csrcs/visibility.c); README "Visibility and spectrum"
% AUTHOR: the observer sees one side; the factor 4 pi in eq. (Lnu) makes the spectrum
%         isotropic-equivalent. Check the wording.

The disc rings radiate as blackbodies at their temperatures~\eqref{eq:Tdisc}, and so does the central emitter of a NS.
The wind rings do not.
Above about 0.3~keV the wind is far too thin to absorb photons and emit them again (Sect.~\ref{sec:disc:thermalization}).
It scatters the disc photons without changing their frequency, so they keep the colour temperature $T_{\rm w,raw}$ of the disc that made them (Eq.~\ref{eq:Twraw}) and are spread over the larger photosphere.
Each wind ring is therefore a diluted blackbody,
\begin{equation}
I_\nu = w\,B_\nu(T_{\rm w,raw}), \qquad w = f_{\rm cor} = \frac{L_{\rm disc,hidden}}{L_{\rm wind,raw}} ,
\label{eq:Inu_wind}
\end{equation}
% code: ULXspectra-mc f546151, csrcs/wind_emission.c (wind_spectrum == 1: f_cor * B_nu(T_j/f_T_cor), f_T_cor = f_cor^(1/4));
%       README "Wind temperature and energy budget" ("Wind spectrum"), "Visibility and spectrum"; PLAN C21
which carries the same energy as a blackbody at $\Tw$ (Eq.~\ref{eq:fcor}).
The isotropic-equivalent spectrum is
\begin{equation}
L_\nu(i) = 4\pi\sum_j I_{\nu,j}\,A_{{\rm proj},j}(i) ,
\label{eq:Lnu}
\end{equation}
% code: isotropic_luminosity_disc_wind() (csrcs/wind_emission.c); README "Visibility and spectrum"
with $I_{\nu,j} = B_\nu(T_j)$ for the disc and the central emitter.
The bolometric luminosities of the wind, the disc and the total follow from the same sum with $\sigma T_j^4/\pi$ in place of $I_{\nu,j}$, where $T_j = \Tw$ for the wind rings.
% code: README "Visibility and spectrum" (Lbol_total, Lbol_wind, Lbol_disc)
For a flat disc without wind, $A_{{\rm proj},j} = \cos i$ times the area of one side of the ring, and the isotropic-equivalent luminosity is $L_{\rm iso} = 2\cos i\,L_{\rm bol}$.
% source: old manuscript (SEA_spectrum.zip, main.tex, "Observed flux"); follows from Eqs. (muproj), (Lnu) with dz/dR = 0
% AUTHOR: blackbody without colour correction or Comptonization (see Discussion);
%         say so here explicitly?

The diluted wind is harder than a blackbody at $\Tw$: its colour temperature is $T_{\rm w,raw} = w^{-1/4}\,\Tw$, and its Rayleigh--Jeans tail is fainter by the factor $w^{3/4}$.
At the 37 old BH best fits with a radiating wind, $w = 0.03$--$0.92$ (median 0.14), so the colour temperature of the wind is 1.04--2.4 times $\Tw$ (median 1.6).
% source: ULXspectra-mc PLAN.md C21 (fits from before the temperature step at R_sph, eq. Tdisc); README "Wind spectrum"; commit f546151
The blackbody at $\Tw$ is the limit of full thermalisation in the wind; it was the rule of earlier code versions and is kept as an alternative (\code{wind\_spectrum} = 0).
% code: wind_spectrum switch (default 1), CHANGELOG 2026-11-04
Since scattering cannot raise the intensity above $B_\nu(T_{\rm w,raw})$, a model with $w > 1$ is flagged (Sect.~\ref{sec:wtemp:budget}); of the 87 old full best fits recomputed with the present rules, only one, NGC300 BH3a, has $w > 1$ (1.03).
% source: CHANGELOG 2026-11-04 ("f_cor > 1 only at NGC300 BH3a (1.03)"); README "Wind spectrum"
The dilution has limits.
In the UV, near 0.01~keV, the wind does thermalise, and there a blackbody at $\Tw$ would be the better description (Sect.~\ref{sec:disc:thermalization}).
One $w$ holds for the whole wind, and each wind ring has one source radius $\xeff$.
% source: README "Wind spectrum" (limitations)

\subsection{Transfer through the wind}
\label{sec:spectrum:transfer}
% Sources: ULXspectra-mc docs/transfer_correction_results.md (commit 2272e8a; state of the model 2f22b12,
%          branch mc-transfer-corrections); README "Visibility and spectrum"; PLAN.md C18

The ray cast alone treats the wind above the surface $\taup = 1$ as transparent, and the light that the surface occults is lost.
A Monte Carlo of electron scattering through the wind shows that this makes $L_{\rm iso}$ too low inside the funnel, by a factor of 1.5--2.4, and too high and too hard beyond its edge, where the hot inner disc stays visible (Appendix~\ref{app:mc}).
% source: ULXspectra docs/green_surface_monte_carlo.md, "Conclusions" items 1-2; commit 4da7c04
We therefore correct the ray cast with a table of Monte Carlo runs (switch \code{mc\_transfer} = 1).
Each emitting ring $j$ of the wind, the disc or, for a NS, the central emitter contributes at $\mu = \cos i$
\begin{equation}
L_j(\mu) = 4\sigma T_j^4\,A_{{\rm att},j}(\mu) + S_j\,h_j(\mu) .
\label{eq:Lj_transfer}
\end{equation}
% code: csrcs/mc_transfer.c; README "Visibility and spectrum" (effective area A_att + S_j h_j/(4 sigma T^4));
%       transfer_correction_results.md sect. 1
Here $A_{{\rm att},j}$ is the projected area of the ring times ${\rm e}^{-\tau}$ along the ray to the observer through the wind of the model, averaged over 16 azimuths.
For a NS the central body, the surface of the magnetosphere, is opaque: a ray that meets it adds nothing to $A_{{\rm att},j}$.
% code: README "Visibility and spectrum" (16 azimuths, 40 log steps per ray; kappa rho from the tau_perp columns;
%       NS central body opaque, one crossing per ring radius)
The light that the ring does not send out directly,
\begin{equation}
\begin{split}
S_j = {} & 2\sigma T_j^4\,2\pi R\,\Delta R\sqrt{1 + \left(\frac{{\rm d}z}{{\rm d}R}\right)^2} \\
& - \int_0^1 4\sigma T_j^4\,A_{{\rm att},j}(\mu)\,{\rm d}\mu ,
\end{split}
\label{eq:Sj}
\end{equation}
% code: README "Visibility and spectrum" (S_j; 10 mu midpoints; traced on every 4th ring, interpolated in ln R)
leaves with the angular profile $h_j(\mu)$, normalised to $\int_0^1 h_j\,{\rm d}\mu = 1$:
\begin{equation}
h_j(\mu) =
\begin{cases}
(1 - f_j)\,g_{\rm in}(u)/(1 - \cos\theta_{\rm c}), & \mu > \cos\theta_{\rm c}, \\
f_j\,g_{\rm out}(v)/\cos\theta_{\rm c}, & \mu < \cos\theta_{\rm c},
\end{cases}
\label{eq:hj}
\end{equation}
% code: README "Visibility and spectrum" (h_j, u, v, f_j); csrcs/mc_transfer.c
with $u = (1 - \mu)/(1 - \cos\theta_{\rm c})$ and $v = \mu/\cos\theta_{\rm c}$.
The share $f_j$ of this light that leaves beyond $\theta_{\rm c}$ and the profiles $g_{\rm in}$ and $g_{\rm out}$ come from a table of 117 Monte Carlo runs (Appendix~\ref{app:mc:table}).
They depend on the class of the ring (wind or disc, and its position $x/\xsp$; for a NS with the dipole surface, the position on the central emitter), on the optical depth $\tau_j$ of the ring towards the angle $\theta_{\rm c} + \min[10^\circ, (90^\circ - \theta_{\rm c})/2]$, and on $\theta_{\rm c}$.
% code: README "Visibility and spectrum" (classes, tau_j direction, interpolation: ln f linear between theta_c nodes,
%       profiles linear in cos theta_c, in u, v and ln tau); csrcs/mc_transfer_table.h
Energy is conserved by construction, and the scattering is grey and coherent, so each ring keeps its spectral shape.
% source: transfer_correction_results.md sect. 1
Against Monte Carlo runs of best fits not used in the table, the corrected model gives $L_{\rm iso}$ within 6 per cent inside the funnel for the BH fits, and within $-22$ to $+23$ per cent beyond its edge (Appendix~\ref{app:mc:heldout}).
% source: transfer_correction_results.md sect. 4: MC/C 0.94-1.03 inside, 0.78-1.23 beyond (BH, held out); commit 2272e8a
% AUTHOR: are all the fits of the paper run with mc_transfer = 1 (README: default 1, "common to all presets")?

% OUTLINE: example spectra and their dependence on the main parameters (mdot_0, f_in,
%          i, r_out, beta_w), each component shown separately; figure from a new
%          script in docs/scripts/ of the code repository.
